\documentclass[11pt]{article}
\usepackage[a4paper,margin=23mm]{geometry}
\usepackage{amsmath,amssymb,amsthm,mathtools,bm}
\usepackage{booktabs,microtype,xurl,hyperref}
\hypersetup{colorlinks=true,linkcolor=blue,urlcolor=blue,citecolor=blue,
pdftitle={P54 Common Spectrum, Physical Channels, and Nested Elimination}}
\setlength{\emergencystretch}{2em}
\newcommand{\tr}{\operatorname{tr}}
\newcommand{\ran}{\operatorname{ran}}
\newcommand{\Span}{\operatorname{span}}
\newcommand{\dd}{\mathrm d}
\newcommand{\HH}{\mathcal H}
\newtheorem{proposition}{Proposition}
% Generated by verify_p54_spectral_channels.py; do not hand edit.
\newcommand{\ChecksPassed}{153}
\newcommand{\ChecksTotal}{153}
\newcommand{\SingletPoleRows}{
0.01352047 & \ensuremath{2.3534868\times10^{-7}} & 15.947239 \\
0.09971569 & 0.00020922086 & 0.049259875 \\
2.82916586 & 0.07540586 & 0.0035014756 \\
}

\title{Route-F P54: Common Scalar Spectrum,\\Interaction-Defined Channels,\\and Source-Preserving Nested Elimination}
\author{P54PQ-v2 four-dimensional mainline}
\date{10 September 2026}
\begin{document}
\maketitle
\begin{abstract}
We keep the existing action and stationary benchmark fixed and construct
tree scalar-exchange form factors from actual Higgs, Yukawa and vector
vertices. The complete physical scalar space has 295 real dimensions;
33 declared real current slots generate a 21-dimensional invariant
spectral subspace, with 12 visible eigenvalue clusters. This is a
channel-dependent exact reduction, not an additional spatial dimension
or a new particle fit. A precise Gaussian elimination law transports the
kernel, vertices and induced source-source term together. Omitting the
last term loses the entire Higgs-pair scalar-exchange response in the
tested final chart. An initial nonzero $NN$--PQ coupling conjecture is
rejected: the physical PQ tangent is the gauge-singlet phase and has no
linear residue in the declared channels. All \ChecksPassed/\ChecksTotal\
checks pass without recomputing cached Hessians or changing parameters.
\end{abstract}

\section{Scope and unchanged theory}
The frozen action remains P54PQ-v2, with
\begin{equation}
3\,16_F+54_{H,\mathbb R}+126_{H,\mathbb C}
+10_{H,\mathbb C}+1_{H,\mathbb C},
\end{equation}
and only two symmetric UV family matrices $h_{\rm raw},f_{\rm raw}$.
We reuse its 328-real-coordinate polynomial and the stationary benchmark
$(\omega,\sigma,v_s)=(1,0.1265,0.25)$ in program units. The gauge coupling
is read from the historical P2 iteration, with the reference convention
$\mu_U=g_U\omega$. Neither its dimensionful unification scales nor the
benchmark's phenomenological viability are recertified here.

The primary Higgs direction is the \emph{exact tree} light eigenvector.
The bosonic-improved vector is not substituted into a claimed tree or
loop-complete amplitude. All numerical Yukawa residues below are
coefficients with $h_{\rm raw},f_{\rm raw}$ factored out. Fermion masses,
flavor mixing and rates are not fitted. The result is a set of amputated
tree \emph{scalar-exchange} kernels, not a complete gauge-invariant
scattering calculation including vector exchange, contact diagrams,
loops and external-leg matching.

The conceptual motivation is the proposed ``visible energy spectrum''
in Ref.~\cite{chen}. Here visibility is defined by actual interaction
vertices, not an arbitrarily chosen geometrical projection. Different
description layers do not imply different numbers of spatial dimensions.
No mechanical rotation, helical orbit or particle-specific geometry is
introduced.

\section{Physical scalar spectrum and interaction vertices}
\subsection{Remove gauge directions, retain physical zero modes}
Let $H=V^{(2)}(x_0)$ and let $G$ be an orthonormal basis of the rank-33
broken gauge orbit $\{R_Ax_0\}$. At a stationary vacuum, gauge invariance
gives $HG=0$. Choose an orthonormal physical slice $B$ such that
\begin{equation}
 BB^T=P_{\rm phys}=I-GG^T,\qquad B^TB=I_{295},\qquad
 H_p=B^THB.
\end{equation}
The full physical tree spectrum has 290 positive eigenvalues and five
zeros: four real Higgs components and one physical PQ mode. A zero
eigenvalue is not by itself a reason to discard a state. Conversely,
eaten Goldstones are not additional physical scalar particles.

For orthogonal eigenspace projectors $\Pi_\lambda$ of $H_p$,
\begin{equation}
 H_p=\sum_\lambda\lambda\Pi_\lambda,\qquad
 \sum_\lambda\Pi_\lambda=I_{295}.
\end{equation}
For canonical free propagation about this vacuum, the corresponding
energy branches are $E_\lambda(\mathbf p)=\sqrt{|\mathbf p|^2+\lambda}$
in natural units. Thus mass is the rest-energy gap of a branch, not an
additional coordinate or an observer-selected energy value.
There are 24 distinct numerical tree eigenvalue clusters, including
zero. This count is not the earlier 35-row SM-irrep census: distinct
irreps can share a mass, and gauge directions have now been removed.

\subsection{Higgs-pair source from the same potential}
Let $q_\alpha$, $\alpha=1,\ldots,4$, be the canonical real components of
the unique tree light doublet, and
$\rho=H_{\rm light}^\dagger H_{\rm light}
=\tfrac12\sum_\alpha h_\alpha^2$.
For SM-singlet scalar fluctuations $s_A$, the potential contains
\begin{equation}
 V\supset\frac12s^TKs+s^TJ_\rho\rho,\qquad
 (J_\rho)_A=\frac14\sum_{\alpha=1}^4
 V^{(3)}[e_A,q_\alpha,q_\alpha]. \label{eq:hsource}
\end{equation}
Equivalently $J_\rho=P_{\rm SM\,singlet}V^{(3)}[\cdot,q,q]$ for one
canonical neutral component $q$. We verify this equality by the full
four-component trace, not by selecting a neutral triplet as a detector.
Differentiating the gauge identity twice gives
\begin{equation}
 (R x_0)^T V^{(3)}[\cdot,q,q]+2(Rq)^THq=0.
\end{equation}
Thus the exact tree condition $Hq=0$ eliminates gauge-orbit contamination
of this source. The cubic jets are recovered from content-addressed
Hessian differences; a central difference of the Hessian of a quartic
polynomial is exact algebraically. No new lattice or parameter scan is
performed.

\subsection{Yukawa sources and their selection rules}
In the fixed common complex scalar/spinor convention,
\begin{equation}
 M_F(x)=h_{\rm raw}\otimes\sum_A Z_h^A x_A
       +f_{\rm raw}\otimes\sum_A Z_f^A x_A.
\end{equation}
For the declared spinor components $l,r$, define
\begin{equation}
 g_{A,h}^{lr}=l^T Z_h^A r,\qquad
 g_{A,f}^{lr}=l^T Z_f^A r. \label{eq:yvertex}
\end{equation}
The real and imaginary parts label Hermitian-current components; the
actual family matrices multiply the corresponding bilinears and are
not new independent fit parameters for different channels. The seven
component pairs are $u_Lu^c,d_Ld^c,e_Le^c,\nu_LN,NN,\nu_L\nu_L,u^cd^c$.
Quark pairs use the first matched colour of the common-phase dictionary;
these are not colour-averaged bilinears. Zero components are retained as
explicit selection-rule tests.

For an unbroken generator the full tensor identity is
\begin{equation}
 R_F^T Z^A+Z^A R_F+\sum_B Z^B(R_s)_{BA}=0. \label{eq:ward}
\end{equation}
This is checked for both tensors and all 12 SM generators. For hypercharge
eigenstates $l,r$ with charges $y_l,y_r$, it implies, in the real scalar
coordinates used here,
\begin{equation}
 iR_Y g^{lr}=(y_l+y_r)g^{lr}.
\end{equation}
This defines allowed support independently of the scalar eigenvalues.
In particular $g_h^{NN}=g_h^{\nu_L\nu_L}=0$, whereas the corresponding
$126$ vertices are nonzero. A zero for a selected coloured component
must not be generalized to all colour contractions.

\subsection{Vector-pair sources from the kinetic term}
With $D=\partial+g A_aR_a$ and canonical real fields,
\begin{equation}
 (M_V^2)_{ab}=g^2(R_a x_0)^T(R_b x_0).
\end{equation}
For a complete degenerate vector eigenspace with projector $P_\beta$,
dimension $d_\beta$ and corresponding generators $\widetilde R_b$, the
averaged scalar--vector--vector vertex is
\begin{equation}
 (J_\beta)_A=\frac1{d_\beta}\tr
 \left(P_\beta\frac{\partial M_V^2}{\partial x_A}\right)
 =\frac{2g^2}{d_\beta}\sum_{b\in\beta}
 e_A^T\widetilde R_b^T\widetilde R_bx_0. \label{eq:vv}
\end{equation}
It multiplies the metric contraction of the two external vector
polarizations; it is not yet a decay rate. The four positive mass
eigenspaces have real multiplicities $8,1,12,12$ and mass squares
$0.00506506,0.02532532,0.31652201,0.32158708$ in $\omega^2$.
Degeneracy averaging is fixed before the probe is evaluated. Every such
source is an SM singlet. For unbroken massless vectors $R_bx_0=0$, so
this tree single-scalar mass vertex vanishes. An independent finite
difference of the kinetic-term mass trace checks its factor of two.

\section{Common poles, channel residues and a minimal realization}
Collect all physical source columns into $J=B^TJ_{\rm raw}$.
In units $\omega=1$ there are 33 real current slots: one Higgs-pair
source, 28 Yukawa slots and four vector-pair sources. Different current
dimensions are explicitly recorded; the Yukawa spurions are factored
out. At Euclidean $t=p_E^2$ the scalar-exchange response is
\begin{equation}
 F_E(t)=J^T(H_p+tI)^{-1}J
       =\sum_\lambda\frac{R_\lambda}{t+\lambda},\qquad
 R_\lambda=J^T\Pi_\lambda J\succeq0. \label{eq:spectral}
\end{equation}
Minkowski continuation is $t=-p^2-i0$. The pole positions are properties
of the common operator; the residues record which channels can access
each pole. Vanishing residue does not change the spectrum or turn an
unseen state into negative mass. These are finite tree spectral-theorem
statements, not an all-orders spectral representation for gauge-fixed
unphysical fields.

For any real channel vector $c$,
$c^TR_\lambda c=\|\Pi_\lambda Jc\|^2\geq0$. Degenerate eigenvectors may
be relabelled without changing $R_\lambda$. The two independent moment
checks are
\begin{equation}
 \sum_\lambda R_\lambda=J^TJ,\qquad
 \sum_\lambda\lambda R_\lambda=J^TH_pJ. \label{eq:moments}
\end{equation}
The machine-readable artifact retains full cross-channel residues,
not only squared diagonal couplings. Relative phases and interference
cannot in general be reconstructed from diagonal weights alone.

\begin{proposition}[Interaction-generated minimal spectral space]
For finite self-adjoint $H_p$, let
\begin{equation}
 \HH_J=\Span\{\ran(\Pi_\lambda J):\lambda\in\operatorname{spec}H_p\}.
\end{equation}
Then $\HH_J$ is the smallest $H_p$-invariant subspace containing
$\ran J$. Its orthogonal complement is invariant and does not contribute
to $F_E$ for this current set. Restricting $H_p,J$ to $\HH_J$ reproduces
the response exactly at every point of the resolvent set.
\end{proposition}
\begin{proof}
Each generator $\Pi_\lambda Jc$ is an eigenvector or zero, so the span
is invariant, and $\sum_\lambda\Pi_\lambda J=J$ puts $\ran J$ inside it.
Conversely, on a finite spectrum each $\Pi_\lambda$ is a polynomial in
$H_p$ by Lagrange interpolation. Any invariant subspace containing $J$
therefore contains every $\Pi_\lambda J$. Self-adjointness makes its
orthogonal complement invariant. In this decomposition $J$ has zero
component on the complement, proving the resolvent identity.
\end{proof}

For the declared component currents, $\operatorname{rank}J=19$ but
$\dim\HH_J=21$, with 12 visible eigenvalue clusters. Thus a mere
projection onto the instantaneous source range loses dynamically
accessible directions. The 21-dimensional exact realization is verified
against the full 295-dimensional response at five momenta. This is not
a universal claim that P54 has only 21 scalar degrees of freedom:
different probes or special fitted family cancellations change the
observable subspace. In particular the component-current set is not
closed under every possible external colour/isospin choice.

\section{Actual singlet response and a rejected PQ hypothesis}
\begin{center}
\begin{tabular}{rrr}
\toprule
$m^2/\omega^2$ & $R_{\rho\rho}/\omega^2$ & $R_{NN}^{\rm unit\ spurion}$\\
\midrule
\SingletPoleRows
\bottomrule
\end{tabular}
\end{center}
The last column means the sum of the two real unit-spurion residue
components, not a numerical flavor-matrix norm. The same three radial
poles appear in both channels but with very different weights. No
mass-list fitting was used to produce this result.

\subsection{The physical PQ zero mode is not an NN pole here}
An initial hypothesis that the $NN$ channel necessarily sees the PQ
zero mode failed in the first run. The reason follows from support,
not from a numerical tuning. At this vacuum the phase variation of the
$126$ vacuum lies in the broken gauge orbit. The physical PQ tangent is
\begin{equation}
 a=\frac{P_{\rm phys}\,\delta_{\rm PQ}x_0}
 {\|P_{\rm phys}\,\delta_{\rm PQ}x_0\|},\qquad
 (I-P_S)a=0.
\end{equation}
There is no renormalizable $S16_F16_F$ term in the frozen action, hence
\begin{equation}
 Z_h^A=Z_f^A=0\quad(A\in S),\qquad g_{NN}^Ta=0.
\end{equation}
The actual $J_\rho$ and vector-pair vertices are also orthogonal to $a$.
All declared linear scalar residues vanish on this physical mode. The
verifier separately checks its support, Yukawa tensor zero, $Ha=0$
and all current contractions. This does not imply that the axion is
unphysical, absent, globally decoupled or massless after QCD effects.
Other vertices, derivative descriptions and anomalies are outside the
declared tree scalar-exchange current set.

\subsection{Recover the existing nonzero sterile--Higgs coefficient as a spectral moment}
Let $B_s$ embed the three massive radial modes and define
$K_s=B_s^THB_s$, $j_\rho=B_s^TJ_\rho$ and
$G_A=(g_{NN}^TB_s)_A f_{\rm raw}$. Then
\begin{equation}
 \mathcal L\supset-\tfrac12s^TK_s s-s^Tj_\rho\rho
 -\left[\tfrac12N^T(M_R+G_As_A)N+\mathrm{h.c.}\right].
\end{equation}
Solving the scalar equation at leading order gives
$s=-K_s^{-1}j_\rho\rho+\cdots$. In the convention
$\mathcal L\supset C_{HN}^{ij}N_iN_j\rho+\mathrm{h.c.}$, with no
$1/2$ in this operator, define physical-slice columns
$\bar g=B^Tg_{NN}$ and $\bar J_\rho=B^TJ_\rho$. Then
\begin{align}
 C_{HN}&=\frac12G^TK_s^{-1}j_\rho
 =\frac{f_{\rm raw}}2\sum_{\lambda>0}
       \frac{\bar g^T\Pi_\lambda\bar J_\rho}{\lambda},\\
 \omega C_{HN}&=-0.0848692011507\,i\,f_{\rm raw}.
\end{align}
This reproduces the prior same-action coefficient and its common
component phase. The use of $g^T$, not $g^\dagger$, is required for this
holomorphic Majorana operator. Positivity applies instead to the
Hermitian-current residue matrix in Eq.~\eqref{eq:spectral}. No singular
zero-mode inverse is used, and no absolute physical Yukawa rate is claimed.

\section{Nested elimination preserves a triple, not just a mass matrix}
\subsection{Exact Gaussian derivation with external currents}
Let $x=(x_r,x_h)$ and let $\eta$ collect real external current sources.
Define the quadratic Euclidean functional
\begin{equation}
 S_E[x,\eta]=\frac12x^TKx-x^TJ\eta-\frac12\eta^TC\eta.
\end{equation}
At its stationary point, the current response is
$F=C+J^TK^{-1}J$. Solving the heavy equation gives
\begin{equation}
 x_h=K_{hh}^{-1}(J_h\eta-K_{hr}x_r).
\end{equation}
Completing the square then gives the exact transformation
\begin{align}
 K'&=K_{rr}-K_{rh}K_{hh}^{-1}K_{hr},\label{eq:schurk}\\
 J'&=J_r-K_{rh}K_{hh}^{-1}J_h,\label{eq:schurj}\\
 C'&=C+J_h^TK_{hh}^{-1}J_h.\label{eq:schurc}
\end{align}
Subsequent integration of $x_r$ yields
\begin{equation}
 \boxed{C+J^TK^{-1}J=C'+J'^TK'^{-1}J'.} \label{eq:response}
\end{equation}
This is the source-dependent form of the Feshbach--Schur construction
\cite{feshbach,cde}. At complex $p_E^2$ these are analytic transposes,
not Hermitian conjugates of the continued kernel. Hermitian positivity
was established on the real spectral residues before continuation.

\begin{proposition}[Nested consistency]
For nested retained coordinate sets and invertible eliminated blocks,
successive exact transformations of $(K,J,C)$ give the same final triple
as simultaneous elimination. Changing the order of the eliminated
blocks gives the same response on their common domain of definition.
\end{proposition}
\begin{proof}
The heavy linear equations have a unique solution for each retained
field and source. Successive substitutions and a simultaneous solution
therefore give the same heavy stationary value. Substitution into the
same quadratic functional fixes separately the coefficients of
$x_r^2$, $x_r\eta$ and $\eta^2$, which are $K',J',C'$. This proves the
triple identity. Alternatively it follows by block Gaussian elimination
on the augmented matrix
\begin{equation}
 \begin{pmatrix}K&-J\\-J^T&-C\end{pmatrix}.
\end{equation}
At exceptional block singularities one must change the split or use
analytic continuation of the full response, not invert a zero mode by
an undeclared regulator.
\end{proof}

The induced $C'(p^2)$ is a source-source term with no retained-field
propagator, but it is generally \emph{nonlocal} and can contain poles.
It becomes a local Wilson expansion only inside a valid momentum
domain. Calling it a ``contact term'' does not mean it is constant or
dispensable. In particular an eigenvalue of an eliminated block can
produce singularities in individual $K',J',C'$ which cancel in the
complete response; it must not automatically be advertised as a new
physical pole.

\subsection{Actual P54 factorization and negative control}
We reuse the upper 55-scalar physical embedding $B_U$. The final chart
$B_L$ contains the actual eight-real-dimensional $10_H$ bidoublet parent
and the physical PQ tangent. The remaining physical coordinates define
$B_M$, giving the exact split
\begin{equation}
 295=9_L+231_M+55_U.
\end{equation}
This is a coordinate factorization of the same quadratic theory, not
an assertion that the nine-coordinate chart is a local SM EFT.
The test keeps the full momentum dependence, all sources and the
induced source-source functional.

At $p_E^2/\omega^2=0.001,0.03,0.2,-0.02+0.003i,-0.265+0.001i$,
direct elimination, $U$ then $M$, and $M$ then $U$ agree. The full
response also agrees with its pole-residue sum and with passive
noncanonical coordinate rescaling, provided the kinetic metric and
vertices are transformed together. These operations are tests of
representation covariance, not physical rotations of matter.

The final chart has no massive SM-singlet scalar. Selection rules give
$J'_\rho=0$, although $F_{\rho\rho}\ne0$. All its scalar-exchange response
therefore resides in $C'_{\rho\rho}$. Omitting $C'$ loses 100\% of this
response at every tested momentum. For the entire normalized current
matrix, the analogous relative Frobenius error ranges from about
$0.366$ to $1$. That matrix norm depends on the declared current units;
it is not a universal percentage error in a physical cross section.
The single-channel $\rho\rho$ negative control needs no fitted Yukawa
matrix and does not depend on relative current rescaling.

\subsection{What this does not close}
An interacting loop-level elimination must transport higher vertices,
the functional measure and regulated determinant as well. The Gaussian
determinant at this fixed background is source independent; dropping it
is harmless for the present tree quadratic response but not for a
claimed full effective action or finite matching. The exact identities
do not complete the outstanding common gauge/Goldstone/ghost/scalar
hard-minus-EFT calculation. A finite derivative truncation preserves
them only to its declared order and within its validity domain.

\section{Extra spatial dimensions: a separate, constrained comparison}
No extra direction is used above. If one independently postulates a
flat interval $y\in[0,L]$ and a free real five-dimensional scalar,
\begin{equation}
 S=\frac12\int\dd^4x\int_0^L\dd y\,
 \left[(\partial_\mu\Phi)^2-(\partial_y\Phi)^2-M_5^2\Phi^2\right],
\end{equation}
Neumann conditions and the self-adjoint eigenproblem give
\begin{equation}
 -f_n''=\lambda_nf_n,\quad f_n'(0)=f_n'(L)=0,\quad
 f_0=L^{-1/2},\quad f_{n>0}=\sqrt{2/L}\cos(n\pi y/L).
\end{equation}
Expanding $\Phi=\sum_n\phi_n f_n$ and using orthonormality yields
\begin{equation}
 m_n^2=M_5^2+(n\pi/L)^2,\qquad
 \frac{m_2^2-m_0^2}{m_1^2-m_0^2}=4.
\end{equation}
Thus two common parameters predict an entire tower. For a declared
probe profile $w(y)$, residues are also fixed by
$|\int_0^L w(y)f_n(y)\dd y|^2$; $w$ may not be chosen separately for
each particle. Changing boundary conditions, potentials or particle
assignments adds assumptions and must be counted. This conditional
construction is not adopted or fitted, and it does not derive chiral
families, charges or the observed SM spectrum. The plane-wave/normal-mode
energy is not a mechanical helical trajectory.

\section{Reproducibility and next gate}
The verifier is \texttt{code/verify\_p54\_spectral\_channels.py}; its JSON
contains all source columns, complete physical mass eigenvalues, residue
matrices, selection rules, source hashes and content-addressed cache
reads. It never overwrites or recomputes the supplied Hessian cache.
Use the pinned \texttt{code/requirements\_p54\_matching.txt} environment:
\begin{quote}\small\ttfamily
python3 route\_f/code/verify\_p54\_spectral\_channels.py\ {\normalfont\textbackslash}\newline
\hspace*{1em}--cache-dir /path/to/tmp/p54\_full\_doublet\_cw
\end{quote}
The matching card, common-phase dictionary and prior $C_{HN}$ output
are required inputs. \ChecksPassed/\ChecksTotal\ tests pass. This count
is a computational regression ledger; the propositions state separately
the assumptions of the mathematical identities.

The next mainline task is to incorporate the complete source-dependent
functional into the same-scheme finite matching workflow. The unknown
family matrices must remain correlated inputs until a constrained fit
exists. The present scalar tree gate is closed; full finite matching,
fermion-inclusive light-state iteration, physical pole widths and a
global flavor/seesaw fit remain open. The extra-dimensional comparison
and the independent soliton line do not block this sequence.

\begin{thebibliography}{9}
\bibitem{chen} B. Chen, \emph{Energy spectrum for elementary particles}
(2022), six-page manuscript,
\url{https://indico.cern.ch/event/1109513/contributions/4969967/}.
Used as conceptual motivation, not evidence of established extra dimensions.
\bibitem{feshbach} G. Dusson, I. M. Sigal and B. Stamm,
\emph{The Feshbach--Schur map and perturbation theory} (2021),
\url{https://arxiv.org/abs/2105.02058}.
\bibitem{cde} B. Henning, X. Lu and H. Murayama,
\emph{One-loop Matching and Running with Covariant Derivative Expansion}
(2016), \url{https://arxiv.org/abs/1604.01019}.
The source-response algebra above is tree level, not their completed
one-loop matching calculation.
\end{thebibliography}
\end{document}
